% ===== 底本 第 4 章《焦点-焦点奇点与近环面纤维化》；片段 tr05b；源块 B175--B188 =====

% [源 B175]
\subsection{单值化}

% [源 B176]
现在我们反过来考虑：给定拉格朗日环面纤维化 $X^{\circ}\to B\setminus\{b\}$，它在 $B\setminus\{b\}$ 上诱导出整仿射结构。我们能否把 $X^{\circ}$ 延拓为 $B$ 上（可能带奇点）的拉格朗日环面纤维化？等价地，能否把整仿射结构从 $B\setminus\{b\}$ 延拓到 $B$？

% [Tutor Note: B176 是「反向问题」的提法：从 $B\setminus\{b\}$ 上的整仿射结构出发，问能否越过 $b$ 延拓；B177--B180 说明单值化是延拓的障碍与不变量，B181 则给出延拓后得到的节点整仿射结构。]

% [源 B177]
\para{环面丛上的拓扑单值化。}回忆：若 $\mathbb{T}^n\to M\xrightarrow{\pi} B$ 是拓扑环面丛，则有诱导的基本群作用 $\Phi\colon\pi_1 (B,b)\to\mathcal{M} (\mathbb{T}^n)\to\mathrm{Aut} (H_1 (F_b;\mathbb{R}))\cong GL(n,\mathbb{Z})$，称为\emph{拓扑单值化}。

% [Tutor Note: B177 的 $\mathcal{M}(\mathbb{T}^n)$ 是环面 $\mathbb{T}^n$ 的映射类群（mapping class group），源文用 $\mathcal{M}$ 记之；拓扑单值化即 $\pi_1(B,b)$ 经 $\mathcal{M}(\mathbb{T}^n)$ 作用到 $H_1(F_b;\mathbb{R})\cong\mathbb{Z}^n$ 的表示，像落在 $GL(n,\mathbb{Z})$ 中。]

% [源 B178]
\para{仿射单值化。}$B$ 上的拉格朗日环面纤维化等价于格纤维化 $\Lambda\to B$，因此 $\pi_1 (B,b)$ 的单值化作用必须保持格 $\Lambda_b$。于是我们得到单值化作用 $\Psi_B\colon\pi_1 (B,b)\to\mathrm{Aut} (\Lambda_b)\cong\mathrm{Aff} (n,\mathbb{Z})$，其中 $\gamma\mapsto\gamma^{\ast} (TB,\Lambda)\cong I\times (\mathbb{R}^n,\mathbb{Z}^n)\big/\left((0,p)\sim (1,\psi_B [\gamma](p))\right)$。

% [Tutor Note: B178 末尾的商 $\big/\left((0,p)\sim(1,\psi_B[\gamma](p))\right)$ 是平坦 $\mathbb{R}^n$-丛按 $\mathbb{Z}^n$ 的平移与 $\psi_B[\gamma]$ 扭合取商；取 $n=1$ 核对：$\mathrm{Aff}(1,\mathbb{Z})$ 的元素即 $p\mapsto kp+c$（$k=\pm1$，$c\in\mathbb{Z}$）。]

% [源 B179]
\begin{theorem*}
设 $\pi\colon (M,\omega)\to B$ 是拓扑单值化矩阵为 $A$ 的拉格朗日环面纤维化，则相应的仿射单值化矩阵为 $\left(A^{-1}\right)^T$。
\end{theorem*}

% [源 B180]
\begin{theorem*}
辛 $A_1$-奇点的仿射单值化为 $\begin{pmatrix}1&1\\0&1\end{pmatrix}$，而 $A_k$ 的仿射单值化为 $\begin{pmatrix}1&k\\0&1\end{pmatrix}$。
\end{theorem*}

% [Tutor Note: B180 与 B179 相容性核对：取 $n=2$、拓扑单值化 $A=\begin{pmatrix}1&1\\0&1\end{pmatrix}$，则 $\left(A^{-1}\right)^T=\begin{pmatrix}1&0\\-1&1\end{pmatrix}$；它与 B180 中 $A_1$ 情形的 $\begin{pmatrix}1&1\\0&1\end{pmatrix}$ 只差基底翻转（转置共轭），不是矛盾。]

% [源 B181]
\para{节点整仿射结构。}考虑局部坐标卡 $U_1=\mathbb{R}^2\setminus\{(x,y)\vert y=0, x\geq 0\}$ 与 $U_2=(0,\infty)\times (-\varepsilon,\varepsilon)\subseteq\mathbb{R}^2$，其转移映射为
\[
\phi (x,y)=\left\lbrace\begin{array}{ll} \begin{pmatrix}1&-1\\0&1\end{pmatrix}(x,y),&y>0;\\ (x,y),&y<0. \end{array}\right.
\]
于是得到拓扑流形 $B$，它同胚于 $\mathbb{R}^2$，连同 $B\setminus\{0\}$ 上的整仿射结构。

% [Tutor Note: B181 的转移映射在 $y>0$ 时为 $(x,y)\mapsto(x-y,y)$（矩阵 $\begin{pmatrix}1&-1\\0&1\end{pmatrix}$），在 $y<0$ 时为恒等；沿绕原点的环路走一圈，单值化即该矩阵，故 $B\setminus\{0\}$ 上的整仿射结构不能延拓到 $0$，$0$ 就是节点。]

% [源 B182]
对辛流形 $(M,\omega)$ 上哈密顿且有效作用的 $\mathbb{T}^n$，其矩映射 $\mu_{\mathbb{T}^n}$ 的像是凸多面体。这可推广如下：设 $M$ 非紧或带边界，则可将 $\mu$ 分解为 $\mu=\Phi\circ\pi$，其中 $\pi\colon M\to B$ 是轨道空间纤维化，$\Phi\colon B\to\mathbb{R}^n$ 是仿射浸入。Symington [Sym03] 的工作表明，该环面纤维化可由三元组 $(B,\mathcal{A},\mathcal{S})$ 编码，其中 $\mathcal{S}$ 是 $B$ 的分层；该三元组称为\emph{环面底}（toric base），并附有一些相容条件。

% [Tutor Note: B182 的三元组 $(B,\mathcal{A},\mathcal{S})$ 中，$\mathcal{A}$ 是整仿射结构、$\mathcal{S}$ 是分层；源文所说的「相容条件」指 $\mathcal{A}$ 限制在每个 stratum 上是整仿射结构，且沿 stratum 的边界（节点）按 B181 的转移映射贴合。]

% [源 B183]
\begin{theorem*}
对于 $n$ 维的环面底，存在唯一的辛流形 $(M,\omega)$，且其上配备唯一的拉格朗日纤维化。
\end{theorem*}

% [源 B184]
粗略地说，其证明方式是：考虑平凡纤维化 $(B\times\mathbb{T}^n,\omega)\to B$，作边界约化，最终得到光滑流形。

% [源 B185]
\subsection{近环面流形}

% [源 B186]
近环面流形是带有只含椭圆型与焦点-焦点型奇点的拉格朗日环面纤维化的辛流形。在维数 $4$ 时，这类流形很自然地给出近环面底 $(B,\mathcal{A},\mathcal{S})$，连同离散点集 $\{b_i\in B^{top}\}$：在每个 $b_i$ 的局部我们得到节点整仿射流形，并且每个 $b\in B\setminus\{b_i\}$ 局部为环面的。

% [Tutor Note: B186 的 $B^{top}$ 指 $B$ 的顶层（top stratum），与术语表「top stratum = 顶层」一致；$\{b_i\}$ 即近环面底中的节点（nodal points），其局部模型由 B181 给出。]

% [源 B187]
包含映射 $(B\setminus\{b_i\},\mathcal{A})\to (\mathbb{R}^4,\mathcal{A}_0)$ \textbf{不是}仿射浸入，但沿某个特征值切开便可得到仿射浸入。当底非紧或没有边界时，辛流形被唯一确定；否则唯一性不再成立。在这个意义下，我们得到 $4$ 维近环面纤维化在微分同胚意义下的分类。例如，正如 [LS03] 中所讨论的：

% [Tutor Note: B187 源文有两处拼写合并：“affineimmersion” 应作 “affine immersion”、“aneigenvalue” 应作 “an eigenvalue”；另有主谓不一致 “The inclusions … is”。底本如此，照译并报告。]

% [源 B188]
\begin{table}[htbp]
\centering
\small
\begin{tabular}{lll}
\toprule
底流形 & 节点数 & 总空间 \\
\midrule
$\mathbb{D}^2$ & $n$ & $\mathbb{C}P^2\#(n+k-3)\overline{\mathbb{C}P^2}$ 或 $\mathbb{S}^1\times\mathbb{S}^3$，$k=\#\text{顶点}$ \\
$\mathbb{S}^2$ & $24$ & K3 曲面 \\
\bottomrule
\end{tabular}
\tabtitle{3}{底流形与节点数}
\end{table}

% [Tutor Note: 表 3（B188）为 3 列：底流形（Base）/ 节点数（# nodes）/ 总空间（Total space）。第一行中 $n$ 与 $k$ 是两个不同计数：$n$ 为第 2 列的节点数，$k=\#\text{vertices}$ 为底的顶点数，故总空间为 $\mathbb{C}P^2\#(n+k-3)\overline{\mathbb{C}P^2}$；第二行 $\mathbb{S}^2$ 底上有 $24$ 个节点，对应 K3 曲面。]
